\scanpage{432}
\mspage{171}
% archive sheet number: 431

projection \inserted{$\mathrm{pr}_2$} de $X'\times X$ sur $X$, donc c'est aussi un morphisme parfait (puisque $X'\longrightarrow e$ est parfait), donc il en est de même du composé

\[
(\mathrm{pr}_2^X\underset{\raisebox{0pt}{\rule{.35pt}{.5ex}\rule{3.4em}{.35pt}\rule{.35pt}{.5ex}}}{f')\circ p_{X,f}}
=\unsure{\bigl((|\bigl(}\underset{\raisebox{0pt}{\rule{.35pt}{.5ex}\rule{3.3em}{.35pt}\rule{.35pt}{.5ex}}}{\mathrm{pr}_2^X(p_X\unsure{)|\}}}\,\mathrm{pr}_2')
=\beta\,\mathrm{pr}_2'\overset{\text{déf}}{=}\overline f.
\]

Je rappelle que \deleted{\unsure{$q_f$}} la section $i_f$ de $C(f)$ sur $X'$ est définie à partir de la section $i_X$ de $(C_X,\deleted{\unsure{\sigma_X}})$ sur $X$. C'est une Hot-équiv, car $\sigma_X=\mathrm{pr}_1^X p_X$ est parfait à fibres asphériques, donc universellement une Hot-équivalence, donc $\mathrm{pr}_1 p_{X,f}\colon C(f)\longrightarrow X'$ est aussi une Hot-équivalence, et \unsure{ladite|toute} section en est donc une aussi.
